\documentclass[11pt,a4paper]{article}
\usepackage[margin=27mm]{geometry}
\usepackage{longtable,booktabs,array,fvextra}
\usepackage{amsmath,amsthm,amssymb,fontspec,unicode-math,hyperref}
\setmainfont{STIXGeneral}
\setmathfont{STIX Two Math}
\setmonofont{Menlo}
\hypersetup{hidelinks}
\usepackage{newunicodechar}
\newfontfamily\numeralfallback{STIX Two Text}
\newunicodechar{∎}{\ensuremath{\blacksquare}}
\newunicodechar{⁰}{{\numeralfallback\char"2070\relax}}
\newunicodechar{₀}{{\numeralfallback\char"2080\relax}}
\newunicodechar{₁}{{\numeralfallback\char"2081\relax}}
\newunicodechar{₂}{{\numeralfallback\char"2082\relax}}
\newunicodechar{₃}{{\numeralfallback\char"2083\relax}}
\newunicodechar{₄}{{\numeralfallback\char"2084\relax}}
\newunicodechar{⁻}{{\numeralfallback\char"207B\relax}}
\newunicodechar{⁸}{{\numeralfallback\char"2078\relax}}
\newunicodechar{ⁱ}{{\numeralfallback\char"2071\relax}}
\newunicodechar{⁴}{{\numeralfallback\char"2074\relax}}
\newunicodechar{⁷}{{\numeralfallback\char"2077\relax}}
\newunicodechar{⁶}{{\numeralfallback\char"2076\relax}}
\newunicodechar{⁵}{{\numeralfallback\char"2075\relax}}
\newunicodechar{₈}{{\numeralfallback\char"2088\relax}}
\newunicodechar{ᵢ}{\textsubscript{i}}
\newunicodechar{ⱼ}{\textsubscript{j}}
\newunicodechar{ₜ}{\textsubscript{t}}
\providecommand{\tightlist}{\setlength{\itemsep}{0pt}\setlength{\parskip}{0pt}}
\newtheorem{theorem}{Theorem}
\newtheorem{proposition}[theorem]{Proposition}
\newtheorem{lemma}[theorem]{Lemma}
\newtheorem{corollary}[theorem]{Corollary}
\theoremstyle{remark}\newtheorem{remark}[theorem]{Remark}
\newcommand{\gen}[1]{\langle #1\rangle}
\newcommand{\normalgen}[1]{\langle\!\langle #1\rangle\!\rangle}
\newcommand{\Z}{\mathbb Z}
\newcommand{\F}{\mathbb F}
\setlength{\emergencystretch}{2em}

\providecommand{\passthrough}[1]{#1}

\title{Three-generator hypertopes of prescribed two-power order}
\author{Human Compression Project}
\date{}
\begin{document}
\maketitle
\begin{abstract}
We recover the complete three-generator construction of Hou, Feng and
Leemans: its defining presentation, exact order and type, regular hypertope,
local cyclic kernels and quotient structure. Three dihedral attachments to an
affine square quotient give the group. The normal-kernel recognition theorem
identifies the source presentation without a finite-order comparison. The
parity correction in the printed structural display remains explicit.
\end{abstract}

\section{Marked groups and geometric scope}
A marking is a generating list of nonidentity involutions. For a subset $J$
of its labels, write $G_J$ for the subgroup generated by those markers. The
intersection property requires $G_J\cap G_K=G_{J\cap K}$. In rank three the
geometric type used here records all three pair-product orders.
A regular hypertope is a thin, residually connected incidence geometry whose
type-preserving automorphism group is transitive on chambers. The marking
need not be string. Thus its group intersection property alone is insufficient
for the geometric assertion.

We use the complete companion proofs UA.SK and UA.SK.Geometry. The former
constructs a group by separately supported local kernels. The latter also
transfers the typed coset intersection property from a regular-hypertope
base and proves that the constructed involutory marking gives the full
regular-hypertope group. All its hypotheses are checked below. Write $D_m$
for the dihedral group of order $m$, and use $[x,y]=x^{-1}y^{-1}xy$.
The exact source edition is arXiv:1901.09034v1, dated 23 January 2019;
its principal construction is Theorem 3.2. The auxiliary square decompositions
are also supplied by UA.SquareDecomposition, with the earlier attribution
acknowledged in the source's Lemma 3.1.

\section{Affine seeds and the source presentation}
We retain the proved affine square model [UA, UA.Affine, UA.Square and
UA.SquareDecomposition]. In coordinates, with the rightmost transformation
acting first, it is
\[
 a(x,y)=(-x,y),\qquad b(x,y)=(y,x),\qquad c(x,y)=(x,1-y).
\]
These generate $W=\Z^2\rtimes D_8$ with complete presentation
\begin{equation}
 \gen{a,b,c\mid a^2=b^2=c^2=(ab)^4=(bc)^4=(ac)^2=1}.           \label{eq:square}
\end{equation}
For $\delta\geq0$, define
\[
 Q_\delta=\begin{cases}
 W/q\Z^2,&\delta\text{ even},\quad q=2^{(\delta+2)/2},\\
 W/q\gen{(1,1),(1,-1)},&\delta\text{ odd},\quad q=2^{(\delta+1)/2}.
 \end{cases}
\]
It has order $2^{\delta+5}$, pair orders $(4,4,2)$ and all Boolean deletions.
To identify the parity relation exactly, use $[x,y]=x^{-1}y^{-1}xy$.
Direct affine multiplication gives translations
\[
 [(ab)^2,c]=(0,-2),\qquad [(ab)^2,(bc)^2]=(-2,-2).
\]
The normal closure of a translation in $W$ is the lattice spanned by its
signed-permutation orbit. Hence, respectively, adding
\begin{equation}
 [(ab)^2,c]^{2^{\delta/2}}=1\quad\hbox{or}\quad
 [(ab)^2,(bc)^2]^{2^{(\delta-1)/2}}=1                          \label{eq:parity}
\end{equation}
to \eqref{eq:square} gives exactly $Q_\delta$.

Take $n\geq10$, $s,t\geq2$, $\ell\geq1$, and
$\delta=n-s-t-\ell\geq0$. Let $P$ have the three involutory markers,
the pair-power relations
\[
 (ab)^{2^s}=(bc)^{2^t}=(ac)^{2^\ell}=1,
\]
the three relations
\begin{equation}
 [(ab)^4,c]=[a,(bc)^4]=[(ac)^2,b]=1,                           \label{eq:normal}
\end{equation}
and the appropriate parity relation \eqref{eq:parity}. This is the defining
presentation of [HFL19, Theorem 3.2], with renamed markers.

\begin{theorem}[Hou--Feng--Leemans, corrected structural display]\label{cor:hypertope}
The presentation $P$ is the supported extension of $Q_\delta$ by the three
marked dihedral quotients of orders $2^{s+1},2^{t+1},2^{\ell+1}$ onto
its pair subgroups. In particular it has order $2^n$, the prescribed pair
orders, and a regular rank-three hypertope with full type-preserving group
$P$. The subgroups
\[
 K_{ab}=\gen{(ab)^4},\qquad K_{bc}=\gen{(bc)^4},\qquad
 K_{ac}=\gen{(ac)^2}
\]
form its direct-product base kernel, of type
$C_{2^{s-2}}\times C_{2^{t-2}}\times C_{2^{\ell-1}}$.
\end{theorem}
\begin{proof}
The three pair supports are an antichain. Their proper marked subgroups
are the trivial group and the reflection subgroups, so the maps to the
base are injective on every proper marked subgroup. UA.SK therefore
constructs $H$.

All relators of $P$ hold in $H$. Each pair-power kernel word is supported
only in its own coordinate, where the third marker is the identity, proving
\eqref{eq:normal}. The two commutator words in \eqref{eq:parity} vanish in
each local pair coordinate: one of the relevant markers is deleted, and
direct substitution gives the identity. Their base values are the specified
translations. Each local dihedral presentation holds in $P$. Within its
pair, each marker inverts the displayed kernel generator; the third marker
centralizes it by \eqref{eq:normal}. Thus all three $K$'s are normal.
Killing them reduces the presentation precisely to \eqref{eq:square} and
\eqref{eq:parity}. UA.Recognition applies.

Its order formula gives
$2^{\delta+5}2^{s-2}2^{t-2}2^{\ell-1}=2^n$ and its local recognition gives
the exact pair orders. The square base is a regular string polytope; its
maximal-parabolic cosets have the needed flag property by the retained
polytope correspondence. The dihedral local markers are involutions, so
UA.SK.Geometry supplies the claimed hypertope and its full group. This
last use is explicit: intersection property alone does not supply the flag
property of a general coset geometry.
\end{proof}

The quotient description is consequently
\[
 P\cong
 (C_{2^{s-2}}\times C_{2^{t-2}}\times C_{2^{\ell-1}})\mathbin{.}
 \begin{cases}
 (D_{2^{(\delta+4)/2}}\times D_{2^{(\delta+4)/2}})\rtimes C_2,
       &\delta\text{ even},\\
 (D_{2^{(\delta+3)/2}}\times D_{2^{(\delta+3)/2}})\rtimes C_2^2,
       &\delta\text{ odd}.
 \end{cases}
\]
Here the extension dot makes no splitting assertion. These are the corrected
parity rows already recorded in the accepted source contract; the source
interchanges them and repeats the first pair word when naming its second
kernel. The present argument preserves both corrections. The decompositions
of the square quotient are the previously proved UA.SquareDecomposition,
not new group classifications.


\section{Attribution and exact boundary}
The construction and its geometric conclusion are the source result.
The proof uses one normal supported-kernel recognition theorem for all
parameters. Its local obligations are the three dihedral maps, the literal
parity relator, the normality of the three kernel subgroups, and exact
identification of the residual square quotient. Those checks have all been
made; no cardinality estimate for the unrecognized presented group is used.
The affine normal form, both lattice quotients and both structural
decompositions have complete proofs in the companion umbrella. Its
regular-polytope correspondence and typed-Helly geometry interface are
explicit dependencies.

This is the full admitted principal construction, with the source's two
structural corrections. It does not absorb the separate conjectural
nonexistence statement numbered 3.3 into a construction theorem. Any answer
to that question has its own proof and provenance. The present paper makes
no claim about correction priority or another source edition.

\begin{thebibliography}{9}
\bibitem{UA}[UA] \emph{Supported kernels, translation lattices and finite marked quotients},
companion umbrella in this volume.
\bibitem{HFL19}[HFL19] D.-D. Hou, Y.-Q. Feng and D. Leemans,
\emph{Existence of regular 3-hypertopes with $2^n$ chambers},
arXiv:1901.09034v1, 23 January 2019; published in
\emph{Discrete Mathematics} 342 (2019), 1857--1863.
\bibitem{MS02}[MS02] P. McMullen and E. Schulte,
\emph{Abstract Regular Polytopes}, Cambridge University Press, 2002.
\end{thebibliography}
\end{document}
